\chapter{Least Squares}

\section{The Problem}\index{least squares}
Many data sets show a roughly straight-line pattern, but no line passes through every
point. We want a line that comes as close as possible to all the points at once. The usual
measure of closeness is the \newterm{sum of squared errors}. For a line $\hat{y} = b_0 + b_1 x$
and data points $(x_i, y_i)$,
$$S(b_0, b_1) = \sum_{i=1}^{n} (y_i - b_0 - b_1 x_i)^2$$
The \newterm{least squares} line is the one that makes $S$ as small as possible.

Squaring the errors has two effects. It keeps positive and negative errors from cancelling,
and it gives extra weight to large errors. It also makes the problem solvable with calculus.

\section{Solving with Calculus}\index{normal equations}
A minimum of $S$ occurs where both partial derivatives are zero. Taking the partial derivative with
respect to $b_0$ and setting it to zero gives
$$\sum (y_i - b_0 - b_1 x_i) = 0 \quad\Longrightarrow\quad b_0 = \bar{y} - b_1 \bar{x}$$
Taking the partial derivative with respect to $b_1$ and setting it to zero gives
$$\sum x_i (y_i - b_0 - b_1 x_i) = 0$$
Substituting $b_0$ into the second equation and simplifying gives
$$b_1 = \frac{\sum (x_i - \bar{x})(y_i - \bar{y})}{\sum (x_i - \bar{x})^2} = \frac{S_{xy}}{S_{xx}}$$
These are the formulas used in the simple linear regression chapter.

\section{A Worked Example}
Use the same five points as the simple linear regression chapter:
$x = 1, 2, 3, 4, 5$ and $y = 2.1, 3.9, 6.2, 7.8, 10.1$. The sums needed are
$$\sum x_i = 15, \qquad \sum y_i = 30.1, \qquad \sum x_i^2 = 55, \qquad \sum x_i y_i = 110.2$$
so $\bar{x} = 3$, $\bar{y} = 6.02$, $S_{xx} = 55 - 5(3)^2 = 10$, and $S_{xy} = 110.2 - 5(3)(6.02) = 19.9$.
Therefore
$$b_1 = \frac{19.9}{10} = 1.99, \qquad b_0 = 6.02 - 1.99(3) = 0.05$$
This matches the fitted line from the previous chapter.

\section{The Matrix Form}\index{normal equations!matrix form}
Write the model as a matrix equation. Let $X$ be the matrix whose first column is all ones and
whose second column is the $x$ values, and let $\mathbf{b} = (b_0, b_1)$. Then
$\hat{\mathbf{y}} = X\mathbf{b}$. Minimizing $\|\mathbf{y} - X\mathbf{b}\|^2$ gives the \newterm{normal equations}
$$X^T X \mathbf{b} = X^T \mathbf{y}$$
For the example above,
$$X^T X = \begin{pmatrix} 5 & 15 \\ 15 & 55 \end{pmatrix}, \qquad
X^T \mathbf{y} = \begin{pmatrix} 30.1 \\ 110.2 \end{pmatrix}$$
Solving this $2 \times 2$ system gives $\mathbf{b} = (0.05, 1.99)$, the same answer.
The matrix form extends directly to several predictors, which is how multiple regression is computed.

\section{Residuals and Fit}
At the least squares solution, the residuals $e_i = y_i - \hat{y}_i$ are perpendicular to the columns
of $X$. In particular, they sum to zero, because the first column of $X$ is all ones. The minimum value
of the sum of squared errors is
$$SS_{res} = \sum e_i^2 = 0.107 \text{ for the example above}$$

\section{When Least Squares Is Not Enough}
Least squares assumes the errors are roughly the same size everywhere, and that the relationship
is linear in the parameters. Outliers can pull the line strongly, because squaring gives large errors
more weight. Robust methods use other measures of error when this is a concern.

\begin{Exercise}[label=ls1, title=Solve by hand]
Find the least squares line for the points $(0,1), (1,3), (2,4), (3,8)$.
\end{Exercise}
\begin{Answer}[ref=ls1]
$\bar{x} = 1.5$, $\bar{y} = 4$, $S_{xx} = 5$, and $S_{xy} = 11$. So $b_1 = \frac{11}{5} = 2.2$ and
$b_0 = 4 - 2.2(1.5) = 0.7$. The line is $\hat{y} = 0.7 + 2.2x$.
\end{Answer}

\begin{Exercise}[label=ls2, title=Predict and check]
Using the line from the previous exercise, predict $y$ at $x = 4$, and find $SS_{res}$.
\end{Exercise}
\begin{Answer}[ref=ls2]
$\hat{y} = 0.7 + 2.2(4) = 9.5$. For the four points, the predictions are $0.7$, $2.9$, $5.1$, and $7.3$.
The residuals are $0.3$, $0.1$, $-1.1$, and $0.7$, so $SS_{res} = 0.09 + 0.01 + 1.21 + 0.49 = 1.80$.
\end{Answer}

\begin{Exercise}[label=ls3, title=Why square?]
Give one reason to minimize the sum of squared errors instead of the sum of absolute errors.
\end{Exercise}
\begin{Answer}[ref=ls3]
Squared errors are differentiable everywhere, so calculus gives a closed-form solution. (Another
valid answer: squared errors give large errors more weight, which makes the fit sensitive to outliers.)
\end{Answer}
